% Curs 11 - Pagina 3 (13-XII-1994).
% Fragment pentru % Curs 11 - Pagina 3 (13-XII-1994).
% Fragment pentru % Curs 11 - Pagina 3 (13-XII-1994).
% Fragment pentru % Curs 11 - Pagina 3 (13-XII-1994).
% Fragment pentru \input{courses/pages/curs11-pagina3-ro}.
% Fara preambul sau schimbare fortata de pagina.
% Este pastrata succesiunea din manuscris; observatiile editoriale sunt comentarii.

\subsection*{Curs 11 - Pagina 3}

\[
\left.\begin{aligned}
P &\trianglelefteq G\\
Q &\trianglelefteq G
\end{aligned}\right\}
\Longrightarrow [a,b]\in P\cap Q=\{e\}.
\]

\[
|P\cap Q|\mid |P|=p,\qquad
|P\cap Q|\mid |Q|=q
\Longrightarrow |P\cap Q|=1.
\]

\[
\Longrightarrow ab=ba=c,\qquad
\operatorname{ord}(c)=pq
\Longrightarrow G=\langle c\rangle
\Longrightarrow G\cong\mathbb{Z}_{pq}.
\]
% Primele randuri presupun ca P si Q sunt normale in G.
% Justificarea pentru ambele afirmatii nu este dezvoltata integral
% in paginile precedente; nu este adaugata aici.

\bigskip
\noindent $|A|\equiv1\pmod p$,\qquad $|A|=ph+1$ \quad (din Sylow III),

\noindent $|B|\equiv1\pmod q$,\qquad $|B|=qt+1$.

\medskip
\noindent $|A|=q\Longrightarrow ph=q-1$, imposibil.

\bigskip
\noindent\underline{Calculul nr. orbitelor acțiunii unui gr. finit}

\noindent\underline{pe o mulțime finită}

\medskip
\noindent Fie $G$ grup, $|G|=m$, și $A$ o $G$-mulțime finită, $|A|=n$.

\medskip
\noindent Dacă $\{z_1,\ldots,z_q\}\subseteq A$ e transversală pentru relația
 de echivalență pe $A$, $\sim_G$, atunci
\[
A=\operatorname{orb}(z_1)\cup\operatorname{orb}(z_2)
\cup\cdots\cup\operatorname{orb}(z_q).
\]

\noindent Dacă $a\in G$, $\operatorname{Fix}a=\{x\in A\mid ax=x\}\subseteq A$.

\medskip
\noindent\fbox{T. lui Burnside}
\qquad
\(\displaystyle q=\frac{\sum_{a\in G}|\operatorname{Fix}a|}{|G|}.\)

\bigskip
\noindent $G=S_n$,\qquad $A=\{1,2,\ldots,n\}$,
\[
S_n\times A\longrightarrow A,\qquad (\sigma,i)\longmapsto\sigma(i).
\]

\noindent Atunci $q=1$,\quad
$\operatorname{Fix}\sigma=\{i\in A\mid\sigma(i)=i\}$.

\[
\sum_{\sigma\in S_n}|\operatorname{Fix}\sigma|=n!.
\]

\bigskip
\noindent\underline{Obs.}: $x\sim_G y\Longleftrightarrow
(\exists)\,a\in G$ a.î. $ax=y$.

\[
\Downarrow
\]
\[
\operatorname{Stab}_G(y)=a\bigl(\operatorname{Stab}_G(x)\bigr)a^{-1}
\Longrightarrow
|\operatorname{Stab}_G(y)|=|\operatorname{Stab}_G(x)|.
\]
.
% Fara preambul sau schimbare fortata de pagina.
% Este pastrata succesiunea din manuscris; observatiile editoriale sunt comentarii.

\subsection*{Curs 11 - Pagina 3}

\[
\left.\begin{aligned}
P &\trianglelefteq G\\
Q &\trianglelefteq G
\end{aligned}\right\}
\Longrightarrow [a,b]\in P\cap Q=\{e\}.
\]

\[
|P\cap Q|\mid |P|=p,\qquad
|P\cap Q|\mid |Q|=q
\Longrightarrow |P\cap Q|=1.
\]

\[
\Longrightarrow ab=ba=c,\qquad
\operatorname{ord}(c)=pq
\Longrightarrow G=\langle c\rangle
\Longrightarrow G\cong\mathbb{Z}_{pq}.
\]
% Primele randuri presupun ca P si Q sunt normale in G.
% Justificarea pentru ambele afirmatii nu este dezvoltata integral
% in paginile precedente; nu este adaugata aici.

\bigskip
\noindent $|A|\equiv1\pmod p$,\qquad $|A|=ph+1$ \quad (din Sylow III),

\noindent $|B|\equiv1\pmod q$,\qquad $|B|=qt+1$.

\medskip
\noindent $|A|=q\Longrightarrow ph=q-1$, imposibil.

\bigskip
\noindent\underline{Calculul nr. orbitelor acțiunii unui gr. finit}

\noindent\underline{pe o mulțime finită}

\medskip
\noindent Fie $G$ grup, $|G|=m$, și $A$ o $G$-mulțime finită, $|A|=n$.

\medskip
\noindent Dacă $\{z_1,\ldots,z_q\}\subseteq A$ e transversală pentru relația
 de echivalență pe $A$, $\sim_G$, atunci
\[
A=\operatorname{orb}(z_1)\cup\operatorname{orb}(z_2)
\cup\cdots\cup\operatorname{orb}(z_q).
\]

\noindent Dacă $a\in G$, $\operatorname{Fix}a=\{x\in A\mid ax=x\}\subseteq A$.

\medskip
\noindent\fbox{T. lui Burnside}
\qquad
\(\displaystyle q=\frac{\sum_{a\in G}|\operatorname{Fix}a|}{|G|}.\)

\bigskip
\noindent $G=S_n$,\qquad $A=\{1,2,\ldots,n\}$,
\[
S_n\times A\longrightarrow A,\qquad (\sigma,i)\longmapsto\sigma(i).
\]

\noindent Atunci $q=1$,\quad
$\operatorname{Fix}\sigma=\{i\in A\mid\sigma(i)=i\}$.

\[
\sum_{\sigma\in S_n}|\operatorname{Fix}\sigma|=n!.
\]

\bigskip
\noindent\underline{Obs.}: $x\sim_G y\Longleftrightarrow
(\exists)\,a\in G$ a.î. $ax=y$.

\[
\Downarrow
\]
\[
\operatorname{Stab}_G(y)=a\bigl(\operatorname{Stab}_G(x)\bigr)a^{-1}
\Longrightarrow
|\operatorname{Stab}_G(y)|=|\operatorname{Stab}_G(x)|.
\]
.
% Fara preambul sau schimbare fortata de pagina.
% Este pastrata succesiunea din manuscris; observatiile editoriale sunt comentarii.

\subsection*{Curs 11 - Pagina 3}

\[
\left.\begin{aligned}
P &\trianglelefteq G\\
Q &\trianglelefteq G
\end{aligned}\right\}
\Longrightarrow [a,b]\in P\cap Q=\{e\}.
\]

\[
|P\cap Q|\mid |P|=p,\qquad
|P\cap Q|\mid |Q|=q
\Longrightarrow |P\cap Q|=1.
\]

\[
\Longrightarrow ab=ba=c,\qquad
\operatorname{ord}(c)=pq
\Longrightarrow G=\langle c\rangle
\Longrightarrow G\cong\mathbb{Z}_{pq}.
\]
% Primele randuri presupun ca P si Q sunt normale in G.
% Justificarea pentru ambele afirmatii nu este dezvoltata integral
% in paginile precedente; nu este adaugata aici.

\bigskip
\noindent $|A|\equiv1\pmod p$,\qquad $|A|=ph+1$ \quad (din Sylow III),

\noindent $|B|\equiv1\pmod q$,\qquad $|B|=qt+1$.

\medskip
\noindent $|A|=q\Longrightarrow ph=q-1$, imposibil.

\bigskip
\noindent\underline{Calculul nr. orbitelor acțiunii unui gr. finit}

\noindent\underline{pe o mulțime finită}

\medskip
\noindent Fie $G$ grup, $|G|=m$, și $A$ o $G$-mulțime finită, $|A|=n$.

\medskip
\noindent Dacă $\{z_1,\ldots,z_q\}\subseteq A$ e transversală pentru relația
 de echivalență pe $A$, $\sim_G$, atunci
\[
A=\operatorname{orb}(z_1)\cup\operatorname{orb}(z_2)
\cup\cdots\cup\operatorname{orb}(z_q).
\]

\noindent Dacă $a\in G$, $\operatorname{Fix}a=\{x\in A\mid ax=x\}\subseteq A$.

\medskip
\noindent\fbox{T. lui Burnside}
\qquad
\(\displaystyle q=\frac{\sum_{a\in G}|\operatorname{Fix}a|}{|G|}.\)

\bigskip
\noindent $G=S_n$,\qquad $A=\{1,2,\ldots,n\}$,
\[
S_n\times A\longrightarrow A,\qquad (\sigma,i)\longmapsto\sigma(i).
\]

\noindent Atunci $q=1$,\quad
$\operatorname{Fix}\sigma=\{i\in A\mid\sigma(i)=i\}$.

\[
\sum_{\sigma\in S_n}|\operatorname{Fix}\sigma|=n!.
\]

\bigskip
\noindent\underline{Obs.}: $x\sim_G y\Longleftrightarrow
(\exists)\,a\in G$ a.î. $ax=y$.

\[
\Downarrow
\]
\[
\operatorname{Stab}_G(y)=a\bigl(\operatorname{Stab}_G(x)\bigr)a^{-1}
\Longrightarrow
|\operatorname{Stab}_G(y)|=|\operatorname{Stab}_G(x)|.
\]
.
% Fara preambul sau schimbare fortata de pagina.
% Este pastrata succesiunea din manuscris; observatiile editoriale sunt comentarii.

\subsection*{Curs 11 - Pagina 3}

\[
\left.\begin{aligned}
P &\trianglelefteq G\\
Q &\trianglelefteq G
\end{aligned}\right\}
\Longrightarrow [a,b]\in P\cap Q=\{e\}.
\]

\[
|P\cap Q|\mid |P|=p,\qquad
|P\cap Q|\mid |Q|=q
\Longrightarrow |P\cap Q|=1.
\]

\[
\Longrightarrow ab=ba=c,\qquad
\operatorname{ord}(c)=pq
\Longrightarrow G=\langle c\rangle
\Longrightarrow G\cong\mathbb{Z}_{pq}.
\]
% Primele randuri presupun ca P si Q sunt normale in G.
% Justificarea pentru ambele afirmatii nu este dezvoltata integral
% in paginile precedente; nu este adaugata aici.

\bigskip
\noindent $|A|\equiv1\pmod p$,\qquad $|A|=ph+1$ \quad (din Sylow III),

\noindent $|B|\equiv1\pmod q$,\qquad $|B|=qt+1$.

\medskip
\noindent $|A|=q\Longrightarrow ph=q-1$, imposibil.

\bigskip
\noindent\underline{Calculul nr. orbitelor acțiunii unui gr. finit}

\noindent\underline{pe o mulțime finită}

\medskip
\noindent Fie $G$ grup, $|G|=m$, și $A$ o $G$-mulțime finită, $|A|=n$.

\medskip
\noindent Dacă $\{z_1,\ldots,z_q\}\subseteq A$ e transversală pentru relația
 de echivalență pe $A$, $\sim_G$, atunci
\[
A=\operatorname{orb}(z_1)\cup\operatorname{orb}(z_2)
\cup\cdots\cup\operatorname{orb}(z_q).
\]

\noindent Dacă $a\in G$, $\operatorname{Fix}a=\{x\in A\mid ax=x\}\subseteq A$.

\medskip
\noindent\fbox{T. lui Burnside}
\qquad
\(\displaystyle q=\frac{\sum_{a\in G}|\operatorname{Fix}a|}{|G|}.\)

\bigskip
\noindent $G=S_n$,\qquad $A=\{1,2,\ldots,n\}$,
\[
S_n\times A\longrightarrow A,\qquad (\sigma,i)\longmapsto\sigma(i).
\]

\noindent Atunci $q=1$,\quad
$\operatorname{Fix}\sigma=\{i\in A\mid\sigma(i)=i\}$.

\[
\sum_{\sigma\in S_n}|\operatorname{Fix}\sigma|=n!.
\]

\bigskip
\noindent\underline{Obs.}: $x\sim_G y\Longleftrightarrow
(\exists)\,a\in G$ a.î. $ax=y$.

\[
\Downarrow
\]
\[
\operatorname{Stab}_G(y)=a\bigl(\operatorname{Stab}_G(x)\bigr)a^{-1}
\Longrightarrow
|\operatorname{Stab}_G(y)|=|\operatorname{Stab}_G(x)|.
\]
